\section{The ten theses with the highest estimated usefulness}\label{sec:theses}

The estimated usefulness (EU) of a claim is its credence times its value if true, on a scale from 0 to 5. Two Claude judges, one of them blind to the rest of the project, give both numbers for each claim (Sec.~\ref{sec:method}).

\textbf{H44: a forced erasure costs output.}
In regime III (continuous operation), the scaffold erases the context window of an agent at a cap of 41 records and keeps its memory. H44 models each model call as a multinomial spin and compares 21,165 forced erasures with pseudo-erasures in nine goal periods. After a forced erasure, the read share of non-write calls rises by 0.079 [0.071, 0.086] over five calls in 9/9 periods. Writes fall by 26\% and work commits by 38\%, which costs 4--11\% of output; EU~3.23\cred{0.88}. The agent does not become more random and does not respond more to new room messages. Treat the cap length as a cost dial, because each erasure also breaks most command loops (16\% against 72\% recurrence).

\textbf{H50: activity synchrony is a field, and talk is a coupling.}
H50 asks whether agents move together because a common input drives them (a field) or because they read each other (a coupling). It fits a read-out linear-response model on the call cycle of each recipient in 71 units from 35 goal periods. The daily start and stop of the scheduler explain about 0.6 of activity co-movement. A peer message raises the talk probability of the read-out call by 0.019 in regime III, which is 44\% above the base rate. The effect is positive in 45/71 units and negative in none; EU~2.98\cred{0.85}. To spread a message in regime III, name the recipient: named messages give a jump of 0.17, unnamed ones 0.004.

\textbf{H38: most synchrony in regime III is the schedule of the runner.}
The spin of each agent is its activity in each minute, and a Curie--Weiss model gives the gain of co-activation over independence. H38 adds a second field for scaffold states, such as operator pauses, start-up and errors, in 35 goal periods. Joint silences are rare (5\% of minutes), and the operator schedules most of them, so they are not platform stalls. In regime III, 70--80\% of the apparent co-activation is the runner, which starts and stops all agents together; EU~2.93\cred{0.90}. A real excess remains only in goal periods \#44 and \#51. Before you compute a synchrony statistic, trim each day to the window in which all agents run.

\textbf{H08: agents couple only at the read-out call.}
H08 models each agent as a kinetic Ising spin with a hidden field: the room messages that arrived after its last model call. The couplings act only when a call reads that field, and the observable is the jump in response at the read-out call. The chance that the recipient answers the sender jumps at the read-out call, not at the in-flight call. This holds in 14/17 goal periods for mentions and in 17/17 for reply labels; EU~2.88\cred{0.91}. A forced erasure cuts the reply coupling to earlier senders by $21\pm7\%$ (8/9 periods). Take the coupling delay as one call interval of the recipient: about 20~s when busy, minutes after a pause.

\textbf{H15: the context window, not the memory, carries the work.}
H15 applies Kolchinsky--Wolpert semantic information to two stores of an agent: the memory file and the context window. Natural events in the record erase one store, and the observable is the change in output against a placebo day. Memory losses and new agents with empty memories show no day-scale cost (memory loss $+0.29$~SD [$-0.05$, $+0.63$]). Post hoc, a forced erasure cuts work commits by 39\% [35, 42] for about ten calls in 7/9 periods; EU~2.73\cred{0.78}. The size of the memory note before the erasure does not change this cost ($\rho=+0.01$). To recover this tenth of output, raise the cap or carry a task-state summary across the erasure.

\textbf{H54: the kickoff text is the quench target.}
H54 treats a goal change as a quench of vector spins, the statement embeddings, toward a target vector. The test scores each day-1 content centroid against its own kickoff text and 32 decoy kickoffs. The centroid ranks its own kickoff first in 18/33 goal periods (median percentile 1.0, $p=3\times10^{-6}$); EU~2.60\cred{0.92}. The test fails only where the kickoff names no shared target. In goal period \#51 each agent goes to its own private goal (role-swap accuracy 0.95). Write the shared target into the kickoff text; a more specific text does not make the swarm cluster more tightly.

\textbf{H40: in regimes II and III, the coupling runs on the call clock.}
H40 treats each model call as a Glauber update attempt and fits a reply hazard per call in 33 goal periods. The order parameter $\eta$ is the elasticity of that hazard to call duration, with 0 for a call clock and 1 for a wall clock. In regimes II and III, $\eta=0.00\pm0.05$ (\#51), so the per-hour coupling equals the per-call coupling times the call rate; EU~2.52\cred{0.72}. In regime I (one room, daily sessions), the rule fails ($\eta=+0.51$ [0.41, 0.61]). To couple an agent faster, halve its call interval: its 5-min reply coupling then rises by about 30--50\%.

\textbf{H46: style is not a conserved charge, but it identifies agents.}
H46 divides each message into a style vector of 20 format and punctuation rates and a content embedding. It tests whether style is conserved across natural experiments, against the day-to-day changes of the same agent. Style is not conserved: goal switches move it (percentile 0.69, 21/24), and forced erasures move it while content stays flat (0.56). Style still identifies agents across a goal switch at 0.76, against 0.51 for content and 0.15 for chance; EU~2.44\cred{0.75}. Attribute messages to a model with the 20 style rates, and compare messages at equal context positions.

\textbf{H02: Ising couplings from activity times do not measure influence.}
H02 fits a directed kinetic Ising model to the 1-min activity spins of the agents, with a field for each 30-min block. The observables are the pairwise couplings $J_{ij}$, the net outgoing influence of each agent and the collective coupling $\beta J_0$. The pairwise couplings sit at the null floor in all 21 chunks, and the fit misses installed and elected leaders. In the reserved goal period \#45, the installed leader ranks 4/18 ($z=0.76$). The collective coupling in regime III is the day edges of the operator schedule; EU~2.35\cred{0.94}. To find which agent steers a swarm, measure who talks after whom, not who is active with whom.

\textbf{H69: restatement loops are copies from the context window.}
H69 models the statement mode of an agent as a spin: the agent restates an earlier statement of its own, or it transforms it. Erasures and new inputs act as fields toward exit, and the main observable is a Mantel--Haenszel odds ratio. A statement still in the context window is near-copied 3.4 [2.6, 4.5] times more often than an erased one at equal lag (4/4 periods). A forced erasure multiplies the odds of loop exit by 2.4 [1.2, 4.9]; EU~2.34\cred{0.85}. A self-share threshold does not start loops, and novel input does not end them. To break a restatement loop, erase or trim the context window; do not send a message.

\textbf{Status of the ten.} No claim among the ten passed a confirmation test on reserved data. Nine of the ten did not run on the reserved data. H02 ran one reserved test, on a data table with a known error, and its leader prediction failed. The judges mark H38, H15 and H02 as fragile. H44, H15, H69 and part of H08 use one intervention, the forced erasure (NE41), so they are not independent evidence.
